% GENERATED FILE. Edit canonical lessons, then run npm run generate:books.
% canonical-source-hash: fnv1a64:821ac9a9524a2ee7
% canonical-lessons: GE-C140-liegen, GE-C140-sich-anziehen, GE-C140-erzaehlen, GE-C140-fuehlen, GE-C140-fallen

\chapter{To lie, To get dressed, To tell, To feel, To fall}
\label{ch:ge-to-lie-to-get-dressed-to-tell-to-feel-to-fall}
\hlchaptermodality{140}

\begin{chapteropening}
\textbf{By the end of this chapter:} \emph{I can say to lie, to get dressed, to tell, to feel and to fall.}

Complete the last lesson of chapter 140: I can say to lie, to get dressed, to tell, to feel and to fall.
\end{chapteropening}

\section[liegen]{liegen — to lie (be lying)}

\label{lesson:GE-C140-liegen}



\begin{quote}
Before the new one: say the German for to lay, then the German for to put.
\end{quote}

\subsection*{You'll want to know: liegen}
\textbf{liegen} — \textquotedblleft{}to lie (be lying)\textquotedblright{}.

\begin{cousinweb}
\textbf{liegen} $\leftarrow$ Old High German \emph{ligen} — inherited; EtymDB 2.1 lists it beside English \emph{lie}, \textquotedblleft{}to lie down\textquotedblright{}, Dutch \emph{liggen} and Gothic \emph{ligan}.
\end{cousinweb}

\begin{grammarlens}[title={What you've built}]
One more word for this chapter.
\end{grammarlens}

\subsection*{Your turn}
\begin{itemize}
\raggedright
  \item \emph{Say it:} \emph{liegen}
  \item \emph{Say it:} \emph{liegen}, once more
  \item \emph{Say it:} say \emph{liegen}
  \item \emph{Recall:} say the German for to lay, then the German for to put, then say \emph{liegen} again
\end{itemize}

\begin{tcolorbox}[breakable,colback=teal!4,colframe=teal!35!black,title={Before you move on}]
What does liegen mean? (\textquotedblleft{}To lie (be lying)\textquotedblright{}.) Say it once more.
\end{tcolorbox}
\section[sich anziehen]{sich anziehen — to get dressed}

\label{lesson:GE-C140-sich-anziehen}



\begin{quote}
Before the new one: say the German for to put, then the German for to lie.
\end{quote}

\subsection*{You'll want to know: sich anziehen}
\textbf{sich anziehen} — \textquotedblleft{}to get dressed\textquotedblright{}.

\begin{cousinweb}
\textbf{anziehen} is \emph{an-}, \textquotedblleft{}on\textquotedblright{}, + the verb \textquotedblleft{}to pull\textquotedblright{}, Old High German \emph{ziohan} $\leftarrow$ Proto-Germanic \emph{*teuhaną} (EtymDB 2.1): to pull your clothes on.
\end{cousinweb}

\begin{grammarlens}[title={What you've built}]
One more word for this chapter.
\end{grammarlens}

\subsection*{Your turn}
\begin{itemize}
\raggedright
  \item \emph{Say it:} \emph{sich anziehen}
  \item \emph{Say it:} \emph{sich anziehen}, once more
  \item \emph{Say it:} say \emph{sich anziehen}
  \item \emph{Recall:} say the German for to put, then the German for to lie, then say \emph{sich anziehen} again
\end{itemize}

\begin{tcolorbox}[breakable,colback=teal!4,colframe=teal!35!black,title={Before you move on}]
What does sich anziehen mean? (\textquotedblleft{}To get dressed\textquotedblright{}.) Say it once more.
\end{tcolorbox}
\section[erzählen]{erzählen — to tell (a story)}

\label{lesson:GE-C140-erzaehlen}



\begin{quote}
Before the new one: say the German for to lie, then the German for to get dressed.
\end{quote}

\subsection*{You'll want to know: erzählen}
\textbf{erzählen} — \textquotedblleft{}to tell (a story)\textquotedblright{}.

\begin{cousinweb}
\textbf{erzählen} is \emph{er-} + \emph{zählen}, \textquotedblleft{}to count\textquotedblright{} $\leftarrow$ Old High German \emph{zellen} $\leftarrow$ Proto-Germanic \emph{*taljaną} (EtymDB 2.1). A story is told out, item by item — English \emph{tell} has both senses too.
\end{cousinweb}

\begin{grammarlens}[title={What you've built}]
One more word for this chapter.
\end{grammarlens}

\subsection*{Your turn}
\begin{itemize}
\raggedright
  \item \emph{Say it:} \emph{erzählen}
  \item \emph{Say it:} \emph{erzählen}, once more
  \item \emph{Say it:} say \emph{erzählen}
  \item \emph{Recall:} say the German for to lie, then the German for to get dressed, then say \emph{erzählen} again
\end{itemize}

\begin{tcolorbox}[breakable,colback=teal!4,colframe=teal!35!black,title={Before you move on}]
What does erzählen mean? (\textquotedblleft{}To tell (a story)\textquotedblright{}.) Say it once more.
\end{tcolorbox}
\section[fühlen]{fühlen — to feel}

\label{lesson:GE-C140-fuehlen}



\begin{quote}
Before the new one: say the German for to get dressed, then the German for to tell.
\end{quote}

\subsection*{You'll want to know: fühlen}
\textbf{fühlen} — \textquotedblleft{}to feel\textquotedblright{}.

\begin{cousinweb}
\textbf{fühlen} $\leftarrow$ Old High German \emph{fuolen} $\leftarrow$ Proto-Germanic \emph{*fōlijaną}, \textquotedblleft{}to feel\textquotedblright{} — inherited, the same verb as English \emph{feel}. (EtymDB 2.1)
\end{cousinweb}

\begin{grammarlens}[title={What you've built}]
One more word for this chapter.
\end{grammarlens}

\subsection*{Your turn}
\begin{itemize}
\raggedright
  \item \emph{Say it:} \emph{fühlen}
  \item \emph{Say it:} \emph{fühlen}, once more
  \item \emph{Say it:} say \emph{fühlen}
  \item \emph{Recall:} say the German for to get dressed, then the German for to tell, then say \emph{fühlen} again
\end{itemize}

\begin{tcolorbox}[breakable,colback=teal!4,colframe=teal!35!black,title={Before you move on}]
What does fühlen mean? (\textquotedblleft{}To feel\textquotedblright{}.) Say it once more.
\end{tcolorbox}
\section[fallen]{fallen — to fall}

\label{lesson:GE-C140-fallen}



\begin{quote}
Before the new one: say the German for to tell, then the German for to feel.
\end{quote}

\subsection*{You'll want to know: fallen}
\textbf{fallen} — \textquotedblleft{}to fall\textquotedblright{}.

\begin{cousinweb}
\textbf{fallen} $\leftarrow$ Old High German \emph{fallan} $\leftarrow$ Proto-Germanic \emph{*fallaną} — inherited, the same verb as English \emph{fall}. (EtymDB 2.1)
\end{cousinweb}

\begin{grammarlens}[title={What you've built}]
One more word for this chapter.
\end{grammarlens}

\subsection*{Your turn}
\begin{itemize}
\raggedright
  \item \emph{Say it:} \emph{fallen}
  \item \emph{Say it:} \emph{fallen}, once more
  \item \emph{Say it:} say \emph{fallen}
  \item \emph{Recall:} say the German for to tell, then the German for to feel, then say \emph{fallen} again
\end{itemize}

\begin{tcolorbox}[breakable,colback=teal!4,colframe=teal!35!black,title={Before you move on}]
What does fallen mean? (\textquotedblleft{}To fall\textquotedblright{}.) Say it once more.
\end{tcolorbox}
